\documentclass[10pt,a4paper]{amsart}
\usepackage[margin=19mm]{geometry}
\usepackage{amsmath,amssymb,mathtools}
\usepackage[scr=rsfs,cal=euler]{mathalfa}
\usepackage{xfrac}
\usepackage[unicode,hidelinks,bookmarks=false]{hyperref}
\newcommand{\ie}{i.e.\ }
\newcommand{\cf}{cf.\ }
\newcommand{\resp}{respectively\ }
\newcommand{\Subsection}{\subsection}
\providecommand{\nobreakdash}{\nobreak\hbox{-}}
\setlength{\emergencystretch}{2em}
\multlinegap0pt
\usepackage{bm}
\usepackage{bbm}
\usepackage{dsfont}
\usepackage{xspace}
\usepackage{cancel}
\usepackage{mathtools}
\usepackage{amsthm}
\makeatletter
\@ifundefined{thm}{\newtheorem{thm}{Thm}}{}
\@ifundefined{lem}{\newtheorem{lem}[thm]{Lemma}}{}
\makeatother
\newcommand{\R}{\mathbb{R}}
\newcommand{\bi}{{\bf i}}
\newcommand{\bj}{{\bf j}}
\newcommand{\bk}{{\bf k}}
\newcommand{\mH}{\mathbb{H}}
\renewcommand{\Im}{\textrm{Im}}
\title[On the sub-Riemannian geometry of the quaternionic Heisenber]{On the sub-Riemannian geometry of the quaternionic Heisenberg group}
\author{Joonhyung Kim, Ioannis D. Platis, Li-Jie Sun}
\date{Source: arXiv:2601.11312v2 (CC BY 4.0)}
\begin{document}
\maketitle

\noindent\emph{Source and licence.} On the sub-Riemannian geometry of the quaternionic Heisenberg group. Version arXiv:2601.11312v2 (CC BY 4.0), distributed under the Creative Commons Attribution 4.0 International licence, as recorded in the arXiv licence metadata for this exact version and shown on the abstract page https://arxiv.org/abs/2601.11312v2. Full rights notice: licenses/arxiv-2601.11312-CC-BY-4.0.txt.\par\smallskip

\noindent\textit{Editorial scope.} Counted unit: the Lem whose source label is \texttt{lem:ver} in the article (source0.tex lines 332--334, source label \texttt{lem:ver}), delivered together with the article's own complete original proof of it (source0.tex lines 335--395, one \texttt{proof} environment of 61 lines). No mathematics is written, repaired or re-derived: statement and proof are the article's own text line for line. Required context retained uncounted: none. Every definition, display and label of those lines is retained unchanged. Omitted from this package and not redistributed: 1--331, 396--997. Nothing inside a retained excerpt is omitted or altered: every retained body is delivered exactly as the article's lines are written, so no omission receipt of any kind accompanies this package. Citations: the retained excerpts carry no citation command at all, so no bibliography entry of the article is transcribed and no reference is redistributed. Reference adaptation: none was needed and none was made; every reference inside the delivered unit resolves inside it, so no unresolved reference or citation is produced. Printed slips kept literal: no printed word, symbol, hypothesis or number of the counted statement or of its proof was altered. Layout and class: the article's own class is \texttt{amsart} with packages including amsmath, stmaryrd, textcomp, cases, amscd, epsfig, graphicx, soul, color, verbatim, amssymb, amsfonts, amsbsy, amsthm; the wrapper is the lane's pinned \texttt{amsart} editorial wrapper at 10pt on A4, which restates the theorem-like environments the retained text needs and adds the article's own macro definitions that the retained text needs (\texttt{\textbackslash Im}, \texttt{\textbackslash R}, \texttt{\textbackslash bi}, \texttt{\textbackslash bj}, \texttt{\textbackslash bk}, \texttt{\textbackslash mH}), with extra wrapper packages (bm, bbm, dsfont, xspace, cancel, mathtools). The delivered numbering is the unit's own. Assets: no retained span contains a figure environment, a graphics-inclusion command, a PDF-page inclusion, a table or a TikZ picture; no image file is copied into this package, the package declares no asset dependency and no image file is redistributed with it. Licence: version arXiv:2601.11312v2 of this article is distributed under the Creative Commons Attribution 4.0 International licence, as recorded in the arXiv licence metadata for this exact version and shown on the abstract page https://arxiv.org/abs/2601.11312v2. A full rights notice is delivered at licenses/arxiv-2601.11312-CC-BY-4.0.txt. Characters: every retained excerpt is pure ASCII, so the delivered engine, XeLaTeX with its default Latin Modern text font, reports no missing character.
\begin{lem}\label{lem:ver}
There are horizontal curves joining the origin with the point $(0, t)\in(\mH, {\emph\Im}\mH).$
\end{lem}

\begin{proof}
Let $ k_i \in \R_*$,  $i = 1, 2, 3, 4$. 
We first consider the curve $\gamma_1$ given by $\gamma_1(\lambda)=(k_1\lambda,0)$, $\lambda\in[0,1]$. We have that $\dot\gamma_1(\lambda)=k_1\xi_1(\gamma_1(\lambda))$ and this curve joins the origin with the point $p_1=(k_1,0)$. We now consider the horizontal curve $\gamma_2$ given by
$$
\gamma_2(\lambda) = \left(k_1 + k_2 \lambda\, \bi,\, -2k_1 k_2 \lambda \,\bi \right), \quad \lambda\in[0,1].
$$
We have that $\dot\gamma_2(\lambda)=k_2\xi_2(\gamma_1(\lambda))$ and $\gamma_2$ joins $p_1$ and the point $ p_2 = \left(k_1 + k_2 \bi,\, -2k_1 k_2 \bi\right) $. 
Similarly, starting from $p_2$, we move in the direction of $k_3\xi_3$ along the curve $\gamma_3$ given by 
$$
\gamma_3(\lambda) = \left(k_1 + k_2 \bi + k_3 \lambda\, \bj,\, -2k_1 k_2 \bi - 2(k_1 k_3 \bj + k_2 k_3 \bk)\lambda \right), \quad \lambda\in[0,1],
$$
which joins $p_2$ and $ p_3 = \left(k_1 + k_2 \bi + k_3 \bj,\ -2(k_1 k_2 \bi + k_1 k_3 \bj + k_2 k_3 \bk)\right) $. Continuing in this fashion, from $p_3$, we move in the direction of the vector field $k_4 \xi_4$ along the curve   $\gamma_4$ given for each $\lambda\in[0,1]$ by
$$
\gamma_4(\lambda) = \left(k_1 + k_2 \bi + k_3 \bj + k_4 \lambda\, \bk,\ -2\left((k_1 k_2 + k_3 k_4 \lambda) \bi + (k_1 k_3 - k_2 k_4 \lambda) \bj + (k_2 k_3 + k_1 k_4 \lambda) \bk\right) \right);
$$ 
this curve joins $p_3$ and the point  
$$
p_4 = \left(k_1 + k_2 \bi + k_3 \bj + k_4 \bk,\ -2\left((k_1 k_2 + k_3 k_4) \bi + (k_1 k_3 - k_2 k_4) \bj + (k_2 k_3 + k_1 k_4) \bk\right) \right).
$$
In a similar manner:
\begin{itemize}
\item
From $p_4$, we travel in the direction of $ -k_1 \xi_1 $, reaching the point  
$$
p_5 = \left(k_2 \bi + k_3 \bj + k_4 \bk,\ -2\left((2k_1 k_2 + k_3 k_4) \bi + (2k_1 k_3 - k_2 k_4) \bj + (k_2 k_3 + 2k_1 k_4) \bk\right) \right);
$$
\item from $p_5$, we follow $-k_2 \xi_2 $, arriving at  
$$
p_6 = \left(k_3 \bj + k_4 \bk,\ -2\left((2k_1 k_2 + k_3 k_4) \bi + (2k_1 k_3 - 2k_2 k_4) \bj + (2k_2 k_3 + 2k_1 k_4) \bk\right) \right);
$$
\item
from $p_6$ we proceed along $ -k_3 \xi_3 $, reaching  
$$
p_7 = \left(k_4 \bk,\ -4\left((k_1 k_2 + k_3 k_4) \bi + (k_1 k_3 - k_2 k_4) \bj + (k_2 k_3 + k_1 k_4) \bk\right) \right)
$$
and finally,
\item from $p_7$, we move in the direction of $ -k_4 \xi_4 $, reaching the point  
$$
p_8 = \left(0,\ -4\left((k_1 k_2 + k_3 k_4) \bi + (k_1 k_3 - k_2 k_4) \bj + (k_2 k_3 + k_1 k_4) \bk\right) \right).
$$
\end{itemize}
Thus, the horizontal path comprising the union of all horizontal paths above joins the origin and the point whose vertical component depends linearly on the quantities $k_1 k_2 + k_3 k_4 $, $ k_1 k_3 - k_2 k_4 $, and $ k_2 k_3 + k_1 k_4 $. By choosing suitable values 
for the coefficients $k_i$, we claim that there always exists a family of horizontal curves connecting the origin to the point $(0, t)$. Actually, if $t=(-4t_1, -4t_2, -4t_3),$ we must solve
\begin{eqnarray*}
k_1 k_2 + k_3 k_4&=&t_1,\\
k_1 k_3 - k_2 k_4&=&t_2,\\
k_2 k_3 + k_1 k_4 &=&t_3.
\end{eqnarray*}
Let $t_2 + i t_3 = R e^{i\phi}.$ Suppose that $z_1 =k_1 + k_2 i,$ and $z_2 = k_3 + k_4i$. Using the second and third equations above, we note that
$$z_1 z_2 = (k_1k_3 - k_2k_4) + i(k_1k_4 + k_2k_3) = t_2 + i t_3.$$
When $R\neq0,$ by setting $z_1 = r_1 e^{i\phi_1}$ and $z_2 = r_2 e^{i\phi_2}$, we derive that $r_1 r_2 = R,$ and $\phi_2 = \phi - \phi_1.$ 
The equation $k_1k_2 + k_3k_4 = t_1$ can be rewritten as 
$$\frac{1}{2} r_1^2 \sin(2\phi_1) + \frac{1}{2} r_2^2 \sin(2\phi_2) = t_1.$$
Substituting the expressions for $r_2$ and $\phi_2$:
$$r_1^2 \sin(2\phi_1) + \frac{R^2}{r_1^2} \sin(2\phi- 2\phi_1) = 2t_1.$$
By letting $X = r_1^2 > 0,$ we get that
$$(\sin 2\phi_1) X^2 - (2t_1) X + R^2 \sin(2\phi - 2\phi_1) = 0.$$
A solution exists if there exists a $\phi_1 \in [0, 2\pi)$ such that the quadratic has at least one positive real root. By the intermediate value theorem, we can select $\phi_1$ such that the coefficients $\sin(2\phi)$ and $R^2 \sin(2\phi - 2\phi_1)$ have opposite signs. In such a case, the product of the roots is negative, ensuring the existence of one positive real root for $X$.
The case where $R=0$ follows trivially by setting $k_1=k_2=0$, yielding $k_3k_4=t_1$, which is solvable for all $t_1$.
 This completes the proof.
\end{proof}

\end{document}
